\documentclass[14pt, a4paper]{article}
\usepackage[a4paper, margin=1in]{geometry}
\usepackage{graphicx} % Required for inserting images
\usepackage{siunitx}

\title{A Comprehensive Dataset on a Continuum Robot Motion}
\author{Andrea Gotelli, Tongjia Zheng, Angela Peloso, Spencer Teetaert, \\Puspita Dewi, Chengnan Shentu, Jessica Burgner-Kahrs}


\newcommand{\ie}{\textit{i.e.}}
\newcommand{\eg}{\textit{e.g.}}

%	Set the bibliography once and for all
\bibliographystyle{ieeetr}

\begin{document}

\maketitle


\begin{abstract}
    Continuum robots are extensively studied with advances in their modeling and an increasing interest in the application of machine learning approaches and the study of their contact and interaction with the environment. 
    However, the continuum robotics fields suffers from the limited availability of continuum robots, and relies on the ad-hoc realization of prototypes.
    Experiments are typically focus on specific tasks or objectives. 
    Although designs and experimental setup are correctly reported in the literature, they are tied to the manufacturing capabilities and sensor available at specific labs. 
    This data descriptor presents a comprehensive and replicable dataset on the motion of a continuum robot with quasi static and dynamics behavior with and without contact with object in the workspace. 
    With the usage of combined sensors measures, we provide a complete dataset with information encompassing both kinematics and force measurements.
\end{abstract}


\section{Background \& Summary}


%   For Data Descriptors, this section should provide an overview of dataset, including the motivation for creating it, as well as outlining the potential reuse value where applicable. 
%   Any previous publications that used these data, in whole or in part, should be cited and briefly summarised. 
%   Please note that because Data Descriptors do not present results or analyses there is no formal mandate to cite a given amount of prior art for comparison, however we recommend that at least some other datasets or outputs relevant to the field are cited for readers' general interest. 
%   Please do not include subjective claims on novelty, impact, or utility.

To provide a comprehensive overview of continuum robotics, this section begins with an overview of Continuum Robots (CR), the precursors to CPRs, which fuelled interest in the field starting in the early 1990s.


CRs are a specific type of robot that is formed with a thin elastic component, which undergoes significant displacements and deformations \cite{Robinson1999} (see Figure \ref{fig:soft_grasping}).
These robots are typically actuated by a special actuation system distributed along its length.
The initial conceptualization of such robots dates back to the 1960s \cite{Anderson1967}, while the late 1990s saw a significant rise in research on continuum robots \cite{Russo2023}, particularly in the fields of manipulation \cite{Kolachalama2020} and medical applications \cite{BurgnerKahrs2015, Dupont2022}.
%
%
Continuum robots are, essentially, manipulators with a design inspired by the bio-kingdom, where the actuation comes from muscles, tendons, fibers and joints. 
%
%
In comparison to rigid link serial robots, these robots exhibit greater flexibility and dexterity. 
The distinction between a robot with rigid links and continuum robots is exemplified in Figure \ref{fig:from_rigid_to_continuum}. 
An example of a continuum robot grasping an object is shown in Figure \ref{fig:soft_grasping}, while other types of CRs are presented in Figure \ref{fig:crs}.
%
%
As they undergo continuous deformation, they exhibit infinite Degrees of Freedom (DoF), \ie{} their shape may undergo changes without affecting the motion of their end. 
Consequently, they can be utilized in situations where the execution of intricate trajectories is required, such as navigating the shape of an object or traversing a curvilinear path.
%
% 
Furthermore, their flexibility has a considerable impact on safety when employed in collaborative tasks, although this does result in a reduction in payload capacity when compared to traditional rigid robotics.


Continuum robot have been studied in their dynamics with different approaches [boyer, renda, shooting]. A new trend is emerging for the study of their collision in quasi static setting [Rao] and dynamics [Matthew].
However, all these methods have been validated with ad-hoc prototype, not easily reproducible and with motion arbitrary chosen for the intended validation.
Moreover, many sensors are related to a specific lab whose purchase and setup has a non negligible economical and time investment.

With this study, we aim to provide a comprehensive dataset of motion for a continuum robot. The design of the robot can be easily replicated and most of the components can be purchased with little economical investment.
Different sensors have been installed on the robot to provide a comprehensive set of measures allowing to benchmark approaches, train data-driven methods and study the interaction with objects in the workspace.
In details, for each of the four cables actuating the robot, tension and pulled length are measured providing information about the internal actuation of the CR. Fiber Bragg Grating Sensor (FBGS) has been used to provide distributed measurements of both the shape of the robot and the internal strain, while a Motion Capture system (OptiTrack) provides absolute measurements of the robot pose, with markers placed on the CR body and on the objects placed in the workspace. 
Additionally, IMU sensors have been installed on the robot body providing orientation and accelerations measures.
Finally, force-torque sensors have been used to measure the wrench at the robot base and the interaction force when contacting an object in the workspace.
The dataset has been structured in separate batches with targeted focus: a first batch study the quasi-static motion of the robot, a second batch study dynamics motions, a third batch introduces gyroscopes on the robot, which slightly degrade the homogeneous inertia of the robot on its arc-lenght, finally a third batch records the robot motion with interaction with an object in the workspace.


A similar dataset has been created for concentric tube robots \cite{Grassmann2022}



\section{Methods}
This section defines the methodology used to collect the data.
We first describe the specifics of each sensor, we then detail how they where mounted on the robot and how we ensured alignment and validity of the recorded data.
Finally, we describe how the data is collected giving a general idea of the data collection pipeline, the programming language and the use of CPU wall-time timestamps. 

\subsection{Design}
Puspita: desribe the design rationale, the constraints given by the usage of the Mark10 sensors and how the setup can be modified (eg. if they don't have the MArk10 sensors).
Detail how the disks are designed and how we can mount sensors on them.

\subsection{Sensors}
In this section, we describe the characteristics of each sensor in terms of frequency readings, range of outputs, and a summary of their specifics available on their datasheet.
\subsubsection{Mark 10}

Force measurements were obtained using a digital force gauge (M3-5, Mark-10, New York, USA). 
The sensor has a measurement range of \SIrange{0}{2.5}{\newton}, a resolution of \SI{0.002}{\newton}, 
and an accuracy of $\pm 0.3\%$ of full scale. The device has an internal sampling rate of 
\SI{2000}{\hertz}. During data collection, the gauge was connected to a PC via USB serial 
communication, and readings were obtained by sending a request command to the device and 
receiving the current force measurement. Due to communication latency and the polling loop, 
the effective acquisition rate was lower than the internal sampling rate and not strictly 
uniform, typically yielding approximately \SIrange{300}{350}{\hertz}.

\subsubsection{FBGS}


Shape measurements were obtained using a fiber Bragg grating (FBG) shape sensing assembly 
(SHAPE A02-1050-N1, FBGS Technologies GmbH, Jena, Germany). The sensing fiber has a 
sensing length of \SI{500}{\milli\meter} with \num{26} gratings distributed along the fiber 
with a pitch of \SI{20}{\milli\meter}. The fiber is embedded in a Nitinol tube with a 
diameter of \SI{380}{\micro\meter} and uses an E2000/PC optical connector. 
The fiber was interrogated using an optical interrogator (FBG-scan 908-EP, FBGS Technologies GmbH), 
configured with \num{8} available channels and a scan rate of \SI{100}{\hertz}. During operation, \num{4} sensing cores were active, 
with integration times of \SIlist{168;99;294;140}{\micro\second} for the respective channels.

\subsubsection{ATI-FT}

Wrench measurements were collected using a 6-axis force/torque sensor 
(MINI40 SI-40-2, ATI Industrial Automation, North Carolina, USA), producing the wrench vector 
$[F_x, F_y, F_z, T_x, T_y, T_z]$. The SI-40-2 calibration provides sensing ranges of 
\SI{40}{\newton} for $F_x,F_y$ and \SI{120}{\newton} for $F_z$, with resolutions of 
\SI{0.01}{\newton} and \SI{0.02}{\newton}, respectively. The torque ranges are 
\SI{2}{\newton\meter} for $T_x,T_y,T_z$ with a resolution of \SI{2.5e-4}{\newton\meter}. 
Analog voltages were acquired using an NI data acquisition device at \SI{1000}{\hertz} per 
channel over a $\pm$\SI{10}{\volt} input range. During logging, \num{5} samples per channel 
were averaged per record and timestamped, resulting in an effective output rate of 
approximately \SI{200}{\hertz}. A bias was estimated from an initial \num{10} samples per 
channel and subtracted from subsequent readings. The calibrated wrench was obtained by 
multiplying the bias-corrected voltage vector by a $6\times6$ calibration matrix.

\subsubsection{Gyroscopes}
Spencer: Add specifics for the Gyroscopes

\subsubsection{Actuators}
The tendons are actuated by two brushless DC motors (Cybergear, Xiaomi, China) with custom 3D-printed tendon spools. They are position controlled at {\color{red}??}\SI{}{\hertz} to track desired tendon displacement or tension inputs, with tension sensing from the Mark 10 sensors for the latter.

\subsubsection{Mocap}
Spencer: Add specifics for the Mocap


\subsection{Experimental setup}
\begin{itemize}
    \item Details on the designed base (cables without pulleys, sturdy, non-occlusion for mocap)
    \item Experiments performed (static loads, planar motions, circle, Lissajous curve) with and without gyroscopes and with interaction with object in workspace.
    \item Experiments starting with planar motion first (alignment of sensors)
    
\end{itemize}

\subsection{Data Collection}
\begin{itemize}
    \item Describe how data is collected
    \item Describe how synchronization is handled (using CPU wall clock)
    \item Note on the rawness of the released data.
\end{itemize}

\section{Data Records}
\begin{itemize}
    \item Data available on GitHub (ContinuumRoboticsLab)
    \item Define structure dataset
    \begin{itemize}
        \item Number of folders
        \item Number of files
        \item Processed and raw data organization
        \item Sensors used and how data is organized for each sensor
    \end{itemize}
\end{itemize}

% \section{Data Overview (optional)}
\section{Technical Validation}
The objective of establishing this dataset of dynamic motion of a continuum robot is the lack of available data from a multi-sensor setup.
Hence, the objective of this section is to ensure that the different measures are synchronized and equally valid. 
We can differentiate the measures in kinematics and force measures. 
The former comprises the OptiTrack Mocap system measuring the absolute poses of the frames placed on the robot, the FBGS measures of strain, the accelerations given by the Gyroscopes and the corresponding reconstructed shape and the motor encoders measuring the rotation angle, which can be converted into a tendon displacement.
On ther other hand, force measurements include the Mark10 measurements of the cable tension, the wrench at the robot base measured by the ATI-FT and the interaction force measured by the Resense FT sensor.

The objective of this section is to show consistency among the different kinematics measure and with the force measurements.
To this aim, we used the reduced-order model of the Geometric Variable Strain \cite{Boyer2021} showing great adaptability and input-output flexibility to address the validity of the recorded dataset.

The Optitrack and FBGS data can be compared by showing consistent behavior of the recorded shape, as represented in Figure X.
By computing numerical derivatives we can obtain an good estimate of the robot acceleration, which is consistent with the IMU measures.

To validate the behavior of the remaining of the dataset, we used the measured tendon tension as input of the GVS-based simulation.
We show that, using the cable tension as simulation input we obtain a consistent behavior of the robot motion and the wrench at the base.
In this part, we would like to draw the attention that the robot parameters have been characterized with some static load. 
To obtain more accurate results, one should identify the parameters of the robot.

To validate the usage of the resense FT sensor to measure the contact forces, we relied on static measures.
We manually apply a force at the robot tip. Knowing the robot shape and the applied force, the corresponding wrench at the base can be computed and compared with the measured values of the ATI-FT.


\section{Usage Notes (optional)}
Add here any issue observed with the data collection

\section{Data Availability}
Many people in Nature Scientific Data use figshare for which I have created an account.
another option could be make a GitHub repository with tha data content (to check if double blind submission).

\section{Code Availability}
Statement about releasing the code for the data acquisition and the postprocessing \& technical validation. 

\section{References}
To be added in the same (standalone) latex file, no \emph{.bib} file included.

\section{Author Contributions}

\section{Competing Interests}

% \section{Acknowledgments (optional)}
\section{Funding}

% \section{Ethics statement}

    

\bibliography{bibliography}

\end{document}
